\documentclass[12pt,letterpaper]{article}
\usepackage[utf8]{inputenc}
\usepackage[spanish]{babel}
\usepackage{amsmath}
\usepackage{amsfonts}
\usepackage{amssymb}
\usepackage[margin=2.5cm]{geometry}

\begin{document}

\begin{center}
    \Large \textbf{Evaluación de Examen 2: Unidad II} \\
    \large Física para Ciencias de la Salud (005-1713) \\
    \vspace{0.5cm}
    \textbf{Estudiante:} Bárbara Rodríguez \quad \textbf{C.I.:} 32.640.760
    \vspace{0.3cm} \\
    \large \textbf{Calificación Final: 5,5/10 puntos}
\end{center}

\vspace{0.5cm}

\section*{Análisis de Resultados}

\subsection*{1. Cinemática (Aceleración media)}
\textbf{Procedimiento y resultado:} La estudiante extrae los datos e identifica la fórmula de aceleración correctamente ($a = \frac{v_f - v_i}{t}$). La operación matemática es correcta ($0,6 / 0,15 = 4$); sin embargo, comete un error de análisis dimensional al reportar el resultado con unidades de velocidad ($4 \text{ m/s}$) en lugar de la unidad correcta para aceleración ($4 \text{ m/s}^2$). \\
\textbf{Veredicto:} \textbf{Parcialmente correcto} (Penalización por unidades incorrectas) (1,5/2 pts).

\subsection*{2. Dinámica (Leyes de Newton)}
\textbf{Procedimiento y resultado:} Identifica correctamente los datos suministrados. No obstante, al momento de operar comete un error conceptual grave: en lugar de aplicar el despeje analítico de la Segunda Ley de Newton para la aceleración ($a = \frac{F}{m}$), plantea erróneamente una multiplicación entre la masa y la fuerza ($25 \text{ kg} \cdot 150 \text{ N}$). Esto arroja un valor incorrecto de $3750$ sin unidades, cuando el resultado real era $6 \text{ m/s}^2$. \\
\textbf{Veredicto:} \textbf{Incorrecto} (Se otorgan puntos parciales por la identificación de datos) (0,5/2 pts).

\subsection*{3. Trabajo Mecánico}
\textbf{Procedimiento y resultado:} Muestra un procedimiento limpio. Plantea correctamente la relación matemática para el trabajo y sustituye los valores ($400 \text{ N} \cdot 0,8 \text{ m}$). El resultado matemático ($320$) es exacto e incluye acertadamente su respectiva unidad del Sistema Internacional en paréntesis (Joules). \\
\textbf{Veredicto:} \textbf{Correcto} (2/2 pts).

\subsection*{4. Energía Potencial}
\textbf{Procedimiento y resultado:} Extrae de manera estructurada los datos (masa, altura y gravedad) y plantea correctamente la fórmula ($E_p = m \cdot g \cdot h$). La multiplicación de los factores es precisa, arrojando la cifra de $17,64$. Sin embargo, omite por completo indicar la unidad de medida (Joules) en el resultado final, incumpliendo la instrucción de expresar todo en unidades del S.I. \\
\textbf{Veredicto:} \textbf{Parcialmente correcto} (Penalización por omisión de unidades) (1,5/2 pts).

\subsection*{5. Potencia Mecánica}
\textbf{Procedimiento y resultado:} La estudiante no presenta ninguna resolución ni intento de desarrollo para este problema en las hojas entregadas. \\
\textbf{Veredicto:} \textbf{Sin responder} (0/2 pts).

\vspace{0.5cm}
\section*{Comentarios Generales y Sugerencias}
El examen denota un buen manejo del orden visual al estructurar los problemas separando primero los datos, pero es imperativo reforzar el aspecto teórico de las unidades y las fórmulas.
\begin{itemize}
    \item Se recomienda repasar los \textbf{despejes algebraicos} correspondientes a las Leyes de Newton. Multiplicar las variables en lugar de dividirlas (como en la Pregunta 2) distorsiona por completo el resultado físico.
    \item Es de suma importancia \textbf{no olvidar las unidades del Sistema Internacional}. Un resultado numérico en física carece de significado clínico o científico si no va acompañado de su magnitud correspondiente (esencial en la pregunta 4 y vital el exponente cuadrado en la aceleración de la pregunta 1).
    \item Se sugiere gestionar mejor el tiempo para evitar dejar problemas totalmente en blanco.
\end{itemize}

\end{document}